\documentclass[a4paper,11pt]{book}
%
% The current document is based on a document by Marten Wortel
% who is greatly acknowledged for making his document available
%
\usepackage{amsmath,amsthm,amsfonts,mathrsfs,anysize,fancyhdr,epsfig}

\newcommand{\bb}{\mathbb}
\newcommand{\mb}{\mathbf}
\newcommand{\mc}{\mathcal}
\newcommand{\pred}{\mathrm{Pred}}

\newtheorem{theorem}{Theorem}
\newtheorem{proposition}{Proposition}
\newtheorem{lemma}{Lemma}
\newtheorem*{corollary}{Corollary}

\theoremstyle{definition}
\newtheorem{definition}{Definition}
\newtheorem{conjecture}{Conjecture}
\newtheorem*{example}{Example}
\newtheorem{algorithm}{Algorithm}


\begin{document}

\begin{titlepage}

\thispagestyle{empty}

{\hspace{10cm}
\begin{minipage}{5cm}
  \includegraphics{tudelft.eps}
\end{minipage}}

\begin{center}

\vspace{1.0cm}

\large
\textbf{\large{Delft University of Technology}} \\
\textbf{\large{Faculty of Electrical Engineering, Mathematics and Computer Science}} \\
\textbf{\large{Delft Institute of Applied Mathematics}}

\vspace{2.0cm}

\framebox[16cm]{\raisebox{0cm}[1.5cm][1.5cm]
 {\textbf{
   \begin{minipage}{13cm}\Large
   \begin{center}
       Considering temporal factors in distributed trust.
   \end{center}
   \end{minipage}}}}

\vspace{2.0cm}

\large A thesis submitted to the\\
Delft Institute of Applied Mathematics\\
in partial fulfillment of the requirements

\vspace{1.0cm}

for the degree\\

\vspace{1.0cm}

\textbf{\large{MASTER OF SCIENCE} \\
in \\
\large{APPLIED MATHEMATICS} \\
\vspace{1cm}
by}

\vspace{1cm}

\textbf{PIM OTTE}\\

\vspace*{\stretch{2}}

\textbf{Delft, the Netherlands \\
August 2016} \\

\vspace{1cm}

Copyright \copyright{} 2016 by Pim Otte. All rights reserved.
\end{center}

\end{titlepage}

\newpage

\thispagestyle{empty}

\quad

\newpage

\thispagestyle{empty}

\hspace{10cm}
\begin{minipage}{5cm}
  \includegraphics{tudelft.eps}
\end{minipage}

\begin{center}
\vspace*{\stretch{1}}

\textbf{\large{MSc THESIS APPLIED MATHEMATICS}}

\vspace{1.5cm}

\textbf{``Considering temporal factors in distributed trust''}

\vspace{1.5cm}

PIM OTTE\\

\vspace{1.5cm}

\textbf{\large{Delft University of Technology}}

\vspace{\stretch{5}}

\end{center}

\begin{center}
\begin{tabular}{ll}
\large{\bf{Daily supervisors}} & \large{\bf{Responsible professor}} \\
$\;$ & \\
\large{Dr.~D.~C.~Gijswijt}  & \large{Prof.\ dr.~ir.~K.~Aardal} \\
$\;$ & \\
\large{Dr.~J.~Pouwelse} & \phantom{niets} \\
\phantom{niets} & \phantom{niets} \\
$\;$ & \\
\large{\textbf{Other thesis committee members\qquad}} & \phantom{niets} \\
$\;$ & \\
\large{Drs.~F.~.I.Ctief } & \large{\dots} \\
$\;$ & \\
\large{\dots} & \large{\dots} \\
$\;$ & \\
\phantom{niets} & \phantom{niets} \\
$\;$ & \\
\large{August 2016} & \large{Delft, the Netherlands}
\end{tabular}
\end{center}
\newpage

\thispagestyle{empty}

\quad

\newpage


\chapter*{\centering \begin{normalsize}Abstract\end{normalsize}}

In this thesis, the problem of estimating trust in a distributed network is considered.

\clearpage
\tableofcontents
\pagenumbering{arabic}
\chapter{Introduction}
    
%Insert general gewauwel here

The internet is more distributed than ever [citation needed].

Distributed systems suffer from a tragedy of the commons.

Defending against sybils is interesting in the sense
that it will also imply that there is no gross unfairness in the system.


\section{Base model}

Traditionally, sybil defense mechanisms are built to exploit some graph,
either based in external trust or interactions. In the case of interactions,
information over time is summarized in a single number. However, this loses
a lot of relevant information.

In order to incorporate temporal information in sybil defense mechanisms,
there will be some assumptions needed about the information which is available.
Hence, applications which fit the following model will be explored.

\begin{definition}[Ordered interaction model]
    An \emph{ordered interaction model} $M$ consists of two sets:  
    
    \begin{itemize}
        \item $P$, a finite set of participants
        \item $I$, a finite partially ordered set of interactions
    \end{itemize}

    Participants represent entities that can interact with each other.
    An interaction consists of two participants, one or both performing
    work for each other. 
    
    An interaction is defined by two pairs of properties:
    \begin{enumerate}
        \item $p, q \in P$, two participants involved in the interaction.
        \item $w_p, w_q \in \bb{R}_{\geq0}$, representing amount of work done by $p$, respectively $q$.
    \end{enumerate}

    Furthermore, for each $p \in P$, the partial order on $I$ must induce a total order on 
    \begin{equation*}
        I_p = \{i \in I : p \mbox{ is a participant in } i\}
    \end{equation*}
\end{definition}

Note that because of the total order on the finite set $I_p$, the following holds:
For each $i \in I_p$, either $i$ is the minimum of $I_p$, or there exists
a unique $i' \in I_p$ such that $i' < i$, and no $i'' \in I_p$ exists such that
$i' < i'' < i$. We define: $\pred(p, i) := i'$, so interaction $i'$ is the predecessor
of interaction $i$ for participant $p$.

  
This model induces a graph that resembles traditional interaction graphs, but preserves the knowledge
of the order of interactions.
\begin{definition}[Ordered interaction graph]
    The directed graph $G_M = (V, A)$ is defined as the \emph{ordered interaction graph} derived from an ordered interaction model $M = (P,I)$.
    Its structure is as follows:

    \begin{itemize}
        \item $V := \{ v_{i} : i \in I\}$\\
        \item $A := \{ (v_{i}, v_{j}) : i, j \in I, \mbox{ s.t. } \exists p \in P, \mbox{ where } i = \pred(p, j)\}$
    \end{itemize}

    In contexts where not defined otherwise, the weight of an arc is the amount of the work performed by $p$ in
    interaction $i$. 
\end{definition}

\subsection{Timeline inference for the base model}

With OIMs, the ordering information of interactions is largely preserved. However, for some
applications it will be useful to be able to associate timestamps with interactions. To
this point, we introduce the concept of timeline inference.

\begin{definition}[Timeline inference for ordered interaction models]
    Let $M = (P, I)$ be an ordered interaction model. 

    $f$ is a \emph{timeline information function} if:
    \begin{itemize}
        \item $J \subseteq I$ 
        \item $f : J \to \bb{R}$
        \item $f$ is increasing with respect to the partial ordering on $I$. That is:
            for $i,j \in J$, if $i \leq j$, then $f(i) \leq f(j)$. 
    \end{itemize}

    By $f[J]$ we denote the range of $f$.
    This allows to define the timeline inference function $T_f$:

    \begin{itemize}
        \item $T_f : I \to f[J] \times f[J]$
        \item $T_f(i) := (\max \{f(j): j \in I, j \leq i\}, \min \{f(j) : j \in I, i \geq j\})$
    \end{itemize}
    
\end{definition}


If $f$ represents some sort of temporal information about the interactions in $J$, 
then the timeline inference function gives interval of minimal size containing valid choices for extrapolating
$f$ beyond elements in $J$. In adddition, for every $j \in J$, $T_f(j) = (f(j), f(j))$.




\section{Base model applied to Tribler}

To obtain more insight into the ordered interaction model, consider the following application:

Tribler is an anonymous peer-to-peer file sharing client. In the application of the model,
the set of participants consists of Tribler clients. Each Tribler client is uniquely identifiable,
but a single real world entity could have multiple participants. In Tribler an interaction
consists of two participants uploading data to each other. The amount of work $w_p$ performed
by a participant $p$ is the amount of megabytes uploaded to their peer ($q$), in MB.

In Tribler, these interactions are recorded in a distributed datastructure, called the multichain.
Every client saves and signs their interactions, marked with a sequence number. They will then request
a digital signature from their peer, who will mark the interaction with their own sequence number.
These sequence numbers induce a partial ordering, where $i < j$, if and only if there exists
a sequence $(i_0, i_1, \ldots i_n)$, where $i=i_0, j=i_n$ and every pair $(i_k, i_{k+1})$ has
a common participant, for which the sequence number in $i_k$ is lower than in $i_{k+1}$.
This results in a full order on $I_p$, since in that case $(i, j)$ is a valid sequence if
$i$ occurred before $j$ and both involve participant $p$.

In Tribler, timeline information inference can be done by each participant $p$, if they
keep a local record of the time at which they participated in interaction $i$. In this
case the timeline information function will be $f : I_p \to \bb{R}$, where $f(i)$ is
the timestamp of interaction $i$.




\end{document}
